\chapter{Experiment}
\label{ch:experiments}

\section{Environment Setup}
\label{sec:env}
We implement the pipeline in Python. The classical and recurrent models run on a CPU; the Transformer
encoders and the QLoRA-fine-tuned LLM need a CUDA GPU, which we run on the MISTI-HPC cluster (an NVIDIA
A40) in \texttt{tmux} sessions. We use the
\texttt{khmer-nltk} segmenter for Khmer word segmentation. Every run writes a self-contained directory
with its configuration, its per-answer predictions for the train/validation/test splits, and a metrics
file, so we can recompute every reported number from saved predictions without re-training. We use a
fixed random seed (42) throughout.

\section{The Experiment Grid}
\label{sec:grid}
The study is organised as a registry of experiments, each isolating a single design choice so that any
change in the result can be traced to one cause. Twelve experiments make up the grid
(Table~\ref{tab:grid}). The first eight (\texttt{v01}--\texttt{v08}) build up the four model pillars
from the cheapest classical baseline to the fine-tuned LLM; the last four (\texttt{exp09}--\texttt{exp12})
are the analyses that turn the trained models into the thesis's findings: the explainability study, the
champion point-metrics, the data-cleaning ablation, and the classical hyperparameter sweep.

Every modelling experiment sweeps the same two text axes from \S\ref{sec:axes} (three preprocessing
modes $\times$ two input formats) crossed with its own model choices, and runs over both dataset
variants (\texttt{full}, \texttt{no10c}). The baseline (\texttt{v01}) and the
BiLSTM grid (\texttt{v05}) are 12 cells per dataset each; the Transformer grid (\texttt{v06}) is 36
cells per dataset (three backbones $\times$ dual/cross $\times$ the text axes). This is why the full
grid runs to hundreds of cells, and why the experiments are split by hardware: the classical and BiLSTM
families run on a CPU, while the Transformer and LLM grids run on the A40 GPU.

\begin{table}[t]
\centering
\caption[The experiment registry]{The experiment registry. Each experiment isolates one design choice
and writes its own leaderboard in the shared schema. ``HPC'' marks the GPU-only experiments.}
\label{tab:grid}
\footnotesize
\setlength{\tabcolsep}{4pt}
\begin{tabular}{@{}llp{8.2cm}@{}}
\toprule
ID & Pillar / role & What it isolates \\
\midrule
\texttt{v01}  & Classical    & TF-IDF cosine and RBF-SVR baseline (12 cells per dataset). \\
\texttt{v02}  & Classical    & Post-hoc threshold calibration: refit the four class cut-points on validation. \\
\texttt{v03}  & Classical    & Add the question max-score as an SVR input feature. \\
\texttt{v03b} & RNN/Transf.  & The max-score feature for the neural heads (BiLSTM CPU, Transformer HPC). \\
\texttt{v04}  & Classical    & A separate TF-IDF+SVR per max-score bucket, routed at test time. \\
\texttt{v05}  & RNN          & Full character BiLSTM + attention grid (12 cells per dataset). HPC. \\
\texttt{v06}  & Transformer  & Full mBERT / XLM-R / GTE dual and cross encoder grid (36 cells). HPC. \\
\texttt{v07}  & Ensemble     & Weighted mean of the top-3 cells across \texttt{v01}--\texttt{v06}. \\
\texttt{v08}  & LLM          & QLoRA fine-tuning of the KhmerGrader family, plus a zero-shot base. HPC. \\
\midrule
\texttt{exp09} & Explainability & SHAP word attribution + plausibility for all four pillars. \\
\texttt{exp10} & Deployment      & Champion point metrics: exact, adjacent, and raw-score error. \\
\texttt{exp11} & Data quality   & The data-cleaning ablation across the two dataset variants. \\
\texttt{exp12} & Tuning         & The classical hyperparameter sweep ($C$ and TF-IDF max features). \\
\bottomrule
\end{tabular}
\end{table}

\section{Before Training Phase}
\label{sec:before}
Before training, we fix the data and the experiment plan above. The classical baseline (\texttt{v01})
runs first because every later experiment is measured against it. The classical enhancements
(\texttt{v02}--\texttt{v04}) each add exactly one idea: calibration refits the class boundaries,
\texttt{v03} hands the model the per-question scale, and \texttt{v04} trains one model per max-score
bucket so that a 10-point and a 20-point question never share a regressor. The neural and Transformer
grids (\texttt{v05}, \texttt{v06}) sweep the same text axes so they are comparable to the classical
cells, and the ensemble (\texttt{v07}) reads the leaderboards of \texttt{v01}--\texttt{v06}, picks the
top three cells by validation QWK, and averages their predictions. The LLM experiment (\texttt{v08}) is
kept deliberately separate from the ensemble pool so that an expensive generative model never silently
inflates the other numbers. Each experiment writes its own leaderboard, and no experiment edits another
experiment's output.

\section{Training Phase}
\label{sec:training}
For every cell, we train on the training split and select on the \emph{validation} split, then report
the chosen configuration \emph{once} on the held-out test split, with no test-set peeking. The neural
and Transformer models use early stopping on validation QWK, which saves compute and guards against
over-fitting on the small corpus.

Figure~\ref{fig:curves} shows the training curves of all three iteratively-trained pillars, the BiLSTM
(left), the GTE encoder (middle), and the fine-tuned LLM (right). In each panel the horizontal axis is
the training epoch and the vertical axis is QWK. For the BiLSTM and the encoder both the training and
validation splits are shown; for the LLM, where evaluating the training split every epoch would require
a full generation pass over it, we track the validation curve only. Two things are visible: the
training and validation lines stay close together where both are shown, so the models are not
over-fitting, and the validation curve flattens after only a few epochs, so there is little to gain from
training longer. Early stopping fires at the epoch where the validation QWK peaks, and that checkpoint,
not the final-epoch one, is what we carry forward to the test set. This is the same selection rule for
every trainable cell. The classical pillar has no such curve: it trains in a single fit, with no epochs.

\begin{figure}[t]
\centering
\includegraphics[width=0.32\textwidth]{bilstm_curve.png}\hfill
\includegraphics[width=0.32\textwidth]{encoder_curve.png}\hfill
\IfFileExists{llm_curve.png}{\includegraphics[width=0.32\textwidth]{llm_curve.png}}%
{\parbox[b]{0.32\textwidth}{\centering\vspace{2.2cm}\emph{LLM validation curve\\(pending re-run)}\vspace{2.2cm}}}
\caption[Training curves of the three trainable pillars]{Training curves by epoch: the BiLSTM (left)
and GTE encoder (middle) show train and validation QWK, while the fine-tuned LLM (right) shows the
validation QWK curve. The lines track each other and plateau early, so early stopping selects the
reported checkpoint without over-fitting. The classical pillar trains in a single fit and has no curve.}
\label{fig:curves}
\end{figure}

We select hyperparameters by validation QWK. For the classical champion, which trains fast enough to
sweep on a CPU, we run a grid over the SVR penalty $C\in\{0.1,1,10\}$ and the TF-IDF maximum features
$\{5,15,30\}\times10^{3}$; the validation-selected configuration changes the \emph{test} QWK by less
than 0.002 versus the defaults (Figure~\ref{fig:hparam}), so the classical result is robust to these
choices. For the neural, Transformer, and LLM pillars, we run the matching learning-rate sweeps
(BiLSTM $\in\{5\mathrm{e}{-}4,1\mathrm{e}{-}3,2\mathrm{e}{-}3\}$; Transformers
$\in\{1\mathrm{e}{-}5,2\mathrm{e}{-}5,5\mathrm{e}{-}5\}$; LLM QLoRA
$\in\{1\mathrm{e}{-}4,2\mathrm{e}{-}4,3\mathrm{e}{-}4\}$) with the same selection rule on GPU hardware.

Figure~\ref{fig:hparam} reads this off directly: for each of the two dataset variants on the
horizontal axis, it places the test QWK of the default configuration next to the test QWK of the
validation-tuned configuration. The two bars are almost the same height in every variant, which is the
point: tuning the classical model's penalty and feature count moves the test QWK by at most a few
thousandths, so the classical result does not depend on a lucky hyperparameter choice.

\begin{figure}[t]
\centering
\includegraphics[width=0.6\textwidth]{fig_hparam_tuning.png}
\caption[Classical hyperparameter sweep]{Classical hyperparameter sweep: default versus
validation-tuned test QWK for each dataset variant. The near-equal bars show the classical model is
robust to its penalty $C$ and TF-IDF feature count.}
\label{fig:hparam}
\end{figure}

\section{Testing Phase}
\label{sec:testing}
In the testing phase we compute every metric from the saved test-set predictions
(\S\ref{sec:metrics}). Each cell writes a self-contained run directory holding a \texttt{config.json}
(exactly what was trained: the model, preprocessing mode, input format, and hyperparameters), a
\texttt{metrics.json} (the train, validation, and test metrics plus the selected epoch), and three
prediction files \texttt{predictions\_\{train,val,test\}.csv} with the per-answer true and predicted
score, label, and raw mark. Every experiment also appends to a shared leaderboard in a common 24-column
schema, so the cells of different experiments line up column-for-column and can be compared directly.

Because every number comes from these saved predictions, we can reproduce every result in
Chapter~\ref{ch:results} without re-training: the entire figure suite is regenerated by a single script
that reads only the CSVs and prediction files (no GPU, no network). We release all code, the dataset and
its two variants, the fixed seed, the per-cell leaderboards, the champion
prediction files, and the figure script. Large model weights (the LLM adapters and encoder checkpoints)
are not tracked in the repository, but are regenerable from the released scripts, and the KhmerGrader
adapters are published separately. A reader with the repository can therefore both recompute every
reported number and retrain any cell from scratch.

\subsection{The Live Teacher-Facing Prototype}
\label{sec:prototype}
We realise the explainability feature as a working web application. The interaction mirrors
what a teacher already does:

\begin{enumerate}[leftmargin=1.8em]
  \item the teacher types (or pastes) the \textbf{question}, the \textbf{reference answer}, and the
  \textbf{student's answer};
  \item the teacher \textbf{chooses a model} from the four pillars;
  \item the system returns a \textbf{score}, a \textbf{word-attribution explanation} (a text-highlight
  heatmap), and \textbf{written feedback}.
\end{enumerate}

The explanation the teacher sees is the same SHAP word attribution that the explainability study
reports, so the explainability study is the feature, not an add-on. The heatmap and feedback always
use the readable segmented tokens (\S\ref{sec:preprocessing}). We
position the tool as a \emph{teacher-assist}, human-in-the-loop system: it supports a teacher's
judgement rather than replacing it. The written feedback comes from an open-source LLM served over an
OpenAI-compatible endpoint, with a rule-based fallback when no model server is reachable, so the tool
never depends on a proprietary grading service. Because SHAP requires many re-scorings per answer, and
for the LLM each one is a full generation, the explanation is offered as a \emph{toggle}: a teacher can
switch it off for a fast score-only result or on for the word-attribution heatmap. The application runs
locally and is not publicly hosted; example screenshots appear in Appendix~\ref{app:prototype}.
